\section{Sujets potentiels}

Les projets ci-dessous illustrent les divers sujets pouvant être abordés dans le cadre du programme. Chaque liste présente le mentor responsable, une référence bibliographique de départ, ainsi qu'un résumé/mots clés du projet.

\begin{projectbox}{Projet 1 -- Méthodes polynomiales en combinatoire}
\textit{Mentor : \textsc{Ismael El Yassini}, étudiant au doctorat de mathématiques -- UQ\`AM}\\
\textit{Livre de référence : Larry Guth, \textit{Polynomial Methods in Combinatorics}, AMS, 2016}

\medskip
\noindent
Ce projet introduit les méthodes polynomiales, qui utilisent des outils algébriques pour résoudre des problèmes combinatoires et géométriques. L’étudiant(e) étudiera l’interpolation polynomiale, les polynômes d’annulation et les arguments de dimension, puis les appliquera au théorème du Nullstellensatz combinatoire, au problème de Kakeya sur les corps finis et au problème des joints. 
\end{projectbox}

\vspace{0.8em}

\begin{projectbox}{Projet 2 -- Introduction à la logique du premier ordre : le théorème de complétude de Gödel}
\textit{Mentor : \textsc{Ahmed Saylani}, \'Etudiant à la maîtrise de mathématiques -- UQ\`AM}\\
\textit{Livre de référence : An Invitation to Mathematical Logic, David Marker}

\medskip
\noindent Ce projet vise à développer les notions syntaxiques et sémantiques nécessaires à l'étude des systèmes formels, puis à établir le lien fondamental entre conséquence logique et démontrabilité.
\end{projectbox}

\vspace{0.8em}

\begin{projectbox}{Projet 3 -- Combinatoire des mots}
\textit{Mentor : \textsc{Yan Lanciault}, \'Etudiant à la maîtrise de mathématiques -- UQ\`AM}\\
\textit{Livre de référence : À déterminer}

\medskip
\noindent
Des études en lettres vous ont toujours intéressé(e), mais vous êtes pris(e) à faire des mathématiques (trop compliquées) ? La combinatoire des mots est là pour vous faire sourire !

\medskip
\noindent
Les notions classiques de la combinatoire des mots seront abordées de manière ludique, motivant l’étude du sujet. La conjugaison, le renversement et l’image de Parikh sont des exemples de notions qui serviront d’appât à la curiosité de l’auditoire. Ensuite, certaines familles de mots classiques, telles que les mots de Fibonacci, de Sturm, de Morse, de Lyndon et de Christoffel, seront introduites afin d’ajouter des personnages intéressants à l’histoire racontée. Finalement, les notions d’automates et d’expressions rationnelles, d’algèbre de Kleene et de langages réguliers seront étudiées. Ceci permettra d’énoncer certaines questions ouvertes intéressantes du domaine et clôturera cette introduction au monoïde libre et à ses dérivés linguistiques.
\end{projectbox}

\vspace{0.8em}

\begin{projectbox}{Projet 4 -- Connexions et courbures en géométrie de faible dimension}
\textit{Mentor : \textsc{Giovanny Ramon Soto Baralt}, étudiant au doctorat de mathématiques -- UQ\`AM}\\
\textit{Livre de référence : Modern Geometry -- Methods and Applications, B.A. Dubrovin, A.T. Fomenko, S.P. Novikov}

\medskip
\noindent Il sera question d'explorer les bases de la géométrie différentielle et riemannienne en dimension deux et trois, d'étudier les équations de courbure, la théorie de jauge et les équations différentielles dans le cadre des variétés.
\end{projectbox}

\vspace{0.8em}

\begin{projectbox}{Projet 5 -- Sujets de base de la géométrie algébrique}
\textit{Mentor : \textsc{Emanuele Ronda}, étudiant au doctorat de mathématiques -- UQ\`AM}\\
\textit{Livre de référence : Undergraduate Algebraic Geometry, M. Reid}

\medskip
\noindent Première approche de la géométrie algébrique : direction à définir avec l'étudiant(e) selon intér\^{e}ts et niveau, par exemple courbes algébriques, lemme de Brody, premières notions d'espaces de modules. 
\end{projectbox}

\vspace{0.8em}

\begin{projectbox}{Projet 6 -- Métriques invariantes sur les espaces homogènes}
\textit{Mentor : \textsc{Pascal Guiffo Kamwa}, étudiant au doctorat de mathématiques -- UQ\`AM}\\
\textit{Livre de référence : An Introduction to Lie Groups and the Geometry of Homogeneous Spaces, Andreas Arvanitoyeorgos}

\medskip
\noindent Dans ce projet, nous étudierons les groupes de Lie et les espaces homogènes d’un point de vue géométrique. Nous nous intéresserons en particulier aux métriques invariantes sur ces espaces et verrons comment les structures algébriques associées aux groupes de Lie permettent de caractériser certaines de leurs propriétés géométriques.
\end{projectbox}

\vspace{0.8em}

\begin{projectbox}{Projet 6 -- Métriques invariantes sur les espaces homogènes}
\textit{Mentor : \textsc{Egor Morozov}, étudiant au doctorat de mathématiques -- UdeM}\\
\textit{Livre de référence : 1. Mauro C. Beltrametti, Ettore Carletti, Dionisio Gallarati, Giacomo Bragadin Monti. Lectures on Curves, Surfaces and Projective Varieties. A classical view of algebraic geometry. EMS Textbooks in Mathematics, European Mathematical Society (EMS), Z\"urich, 2009.}

\medskip
\noindent A basic algebraic result states that a nonzero polynomial of degree $n$ with real coefficients has at most $n$ distinct roots. An upgrade of this result is the ``weak'' B\'ezout's theorem: if two polynomials $P(x,y)$ and $Q(x,y)$ of degrees $m$ and $n$ with real coefficients have finitely many common zeroes, then this number does not exceed $mn$. These results are already very useful, for example, the second result gives a short and elegant proof of Pascal's hexagon theorem. However, one might still feel himself unsatisfied (as some mathematicians of the past did). After all, why do we speak about an upper bound on the number of zeroes rather than the exact number? Developing this idea results in a sequence of corrections to the statements, the most important being the notion of the \emph{intersection multiplicity}. Loosely speaking, the intersection multiplicity is the number of common zeroes of $P$ and $Q$ that arise from a single common zero by a generic perturbation of $P$ and $Q$. Then, to obtain an exact equality in B\'ezout's theorem, one counts each common zero as many times as its multiplicity is. The goal of this project is to develop the theory required to define intersection multiplicity from scratch in a classical geometric style, paying attention to as many details and applications as possible. 
\end{projectbox}